\chapter{はじめに}
\chapterauthor{羽藤英二}

人口減少が進んでいた地方都市を襲った巨大災害直後の破壊の風景の中に立ったとき，何かできたことはなかったか，どうすればよかったのかと思わずにはいられなかった．今，都市と交通を学ぶことにどのような意味があるのだろうか．あるいは激変する社会の中で今私たちに何ができるだろうか．おそらくそれは，道路や鉄道，駅前広場や沿道空間の形を単に学ぶことではない．現実の空間とネットワークの上で，人びとがLLMが得意とする言語のやりとりを超えて「空間」をどのように認知し，選択し，互いに影響し合い，その集積として流れ$=$交通が生まれ，外部不経済としての混雑や滞留，さらには回遊，分断，格差がどのように発生しているのかを捉え直すことである．そして，よりよい理解に基づく行為の反復の果てに，都市空間そのものの価値と構造が歴史的な過程を通じていかにその地理的関係が変容していくのか，徹底的に理解することなくして，いかなる政策も計画も設計もサービスも，威勢のいい掛け声をかけたところで何も変わらないのではないだろうか．

行動モデル，ネットワーク配分，最適化，機械学習という本書の四つの柱は，この過程を，個人の意思決定から社会的な空間変容まで貫いて理解するための基礎理論である．本書の目的は，それらを別々の計算技術として紹介することではない．研究の作業と思考の往復のなかで私たちが考えてきたように，人びとの選択の連鎖がネットワーク現象を生成し，その現象を変えようとしてなされる現実の都市設計や交通計画の行為が，再び人びとの行動そのものを変え，その反復が都市自体を更新していく．この連関を，私たち自身が経験してきた都市設計，交通計画，そしてそれらを下支えする基礎理論の構築と実証の蓄積を踏まえながら，ひとつの連続した知の体系として描き直すことに本書の狙いがある．

その出発点のひとつが，\citet{McFadden1974}によって基礎づけられた離散選択理論である．2000年のノーベル経済学賞へと結実したこの理論は，人はなぜこの経路を選び，この交通手段を選び，この場所へ向かうのかという問いに対し，効用と確率的選択という考え方によって，個人の意思決定を政策や空間設計に応答する構造として厳密に記述する道を開いた．こうして離散選択モデルは世界中の交通計画に利用され，インフラ整備の基礎理論として定着するに至った．ここで重要なのは，ネットワーク上の行動が単なる心理学的記述にとどまらず，制度変更や空間条件の変化に対して地理的分布として応答する対象となったことである．この視点があったからこそ，交通行動は社会基盤計画の中心に置かれ，公共交通計画，都市計画，駅まち設計，自動車交通制御に共通する理論的基礎が築かれたといえよう．

しかし，個人の選択を理解するだけでは都市は見えない．選択はネットワーク上で重なり合い，混雑，所要時間，遅延，容量制約，乗換え負荷，滞在時間といった現象へと変換されるからである．マルコフ過程はこうした行動の連鎖を記述する上で最も基礎的な理論であり，最短経路探索，利用者均衡，確率的配分，動的配分と結びつくことでネットワーク上の行動記述という点で大きな力を発揮することとなった．\citet{Wardrop1952}と\citet{Beckmann1956}は，この問題をネットワーク上の均衡として定式化する道を拓いた．さらに，物流や配車，公共交通運用に目を向ければ，単なる経路選択にとどまらず，誰がどの需要をいつどの順序で担うかという組合せ設計が前景化し，\citet{DantzigRamser1959}以来の VRP，さらには列生成，双対問題，\citet{Benders1962}に代表される分解法が，現実の大規模問題に対抗するための基礎となってきた．こうしたアルゴリズムの実装は，線形代数に基礎づけられた行列計算によって支えられている．そして近年では，ChatGPT の登場によって，その実装や試行錯誤が以前よりはるかに身近なものとなったといってよいだろう．

さらに，現実の交通と都市は，唯一の均衡点に収束する静的な系としてではなく，学習し適応する動的な系として再定義されるべきであろう．制度が変われば人は行動を変え，情報が変われば期待も変わる．\citet{Sandholm2002,Sandholm2010}による人口ゲームと進化動学は，この適応の反復そのものを理論化した点で，従来の均衡配分を決定的に拡張する契機となった．さらにいえば，交通政策とは，ある時点の均衡をずらす一度きりの操作ではない．人びとの模倣，学習，反応，慣性の上に新たな安定状態を形成していく過程への介入なのである．近年の Matsunaga and Hato（2026）は，この系譜を受け継ぎつつ，部分可制御性の下で平均場的インセンティブと逆オークションを結びつけ，都市鉄道混雑制御を設計問題として捉え直している．そこでは，限られた予算，私的情報，日々の適応，平均場的結合が，ひとつの閉ループ系として統合されつつある．

同時に，日々の経路変更や需要変動を実証的に捉えるには，個票レベルの長期観測と，それを扱いうる学習理論が必要になる．Ogawa and Hato（2026）は，1か月規模の個人 GPS 軌跡から日々の経路変更を捉え，Temporal Boltzmann Machine と Graph Attention Network，さらに correlated equilibrium の制約を組み合わせることで，共有情報に対する異質で相関した適応反応を表現しようとしている．ここで重要なのは，深層学習を単なる高精度予測器として用いるのではなく，信念更新や相関した適応の代理構造をもつ実証モデルとして位置づけている点である．学習は理論を置き換えるのではなく，理論が十分に捉えきれてこなかった観測の厚みに踏み込むための手段となる．

この文脈の上に，Transformer や大規模言語モデルを起点とする近年の AI の急速な発展が重なる．\citet{Vaswani2017}以降，系列の長距離依存を捉える注意機構，表現学習，生成モデルは，活動スケジュール，経路系列，都市記述，需要変動の理解に新しい可能性を開いた．さらに，物理法則や保存則を学習に埋め込む PINN は，\citet{Raissi2019}以降，データ駆動型学習と数理モデルの橋渡しとして注目されている．しかし，ここで立ち止まらなければならない．都市と交通の本質は，言語の中だけに存在するのではない．人は文を読んで移動するのではなく，身体をもち，制約を受け，他者と干渉し，空間の抵抗や誘因の中で行動する．そして，その行動は実際にネットワーク負荷を変え，土地利用を変え，拠点の価値を変え，都市そのものを作り替えてしまう．私たちは世界各地の現実の都市や駅を設計し，交通計画を評価し，都市計画に取り組んできた．

その経験から考えるとき，既存の理論はなお，本質的な都市の問題を記述するには十分ではないのではないか．たしかに LLM は強力である．しかしなお最終理論としての弱点がないわけではない．都市と交通に必要なのは，LLM の基盤となる「言語」を超えて，現実の「空間」に作用する物理的な選択，相互作用，制約，学習，設計を統合する理論である．人の行動がネットワーク上でどのように干渉し合い，その反復がどのように都市の形を変え，さらにその変わった都市が次の行動をどのように規定するのか．この循環を捉えるには，効用理論，均衡理論，進化動学，最適化，深層学習，さらには物理制約までを横断しなければならない．都市を「語れる」知能ではなく，都市を「説明し，予測し，設計できる」知能をめぐって世界中の研究者たちが訥々と自分たちの課題と向き合い続けている．

最前線のひとつは自動走行であろう．従来のルールベースを超える E2E 的発想では，周囲の行動を外生的に予測し，その上で自車だけを最適化すればよいと考えられがちであった．しかし実際の交通は，相手もまたこちらを見て応答する相互作用の場である．Furuhashi and Hato（2026）は，この限界を正面から捉え，予測依存型の制御から，相互作用そのものを均衡埋め込みで扱う I2I 型の自動運転へと理論の重心を移している．そこでは，shared constraints の下での variational generalized Nash equilibrium，ICNN による学習可能なコスト設計，そして equilibrium-embedded predictive control が統合され，E2E を越えて現実空間の内生性に向き合う知能の姿が模索されている．さらに彼らは，この枠組みを，これまでにない歩行者と自動走行の混在空間における最適制御と最適設計の理論へと拡張しようとしている．

本書が学部生・大学院生に伝えたいのは，まさにこうしたダイナミックな研究のプロセスとその基礎となる理論のバリエーションである．行動モデル，ネットワーク配分，最適化，機械学習は，互いに隣接する技法ではない．それらは，人間の行動がネットワークを通じて都市を変え，その都市の設計が再び人間の行動を変えるという現実の循環構造に迫ろうとする，ひとつの知の系譜である．\citet{McFadden1974}から\citet{Wardrop1952}，\citet{Beckmann1956}，\citet{Sandholm2010}，\citet{Vaswani2017}，\citet{Raissi2019}へと至る流れは，単なる理論の並列ではなく，観測，均衡，設計，学習を再接続し続けてきた歴史である．そして Matsunaga and Hato（2026），Ogawa and Hato（2026），Furuhashi and Hato（2026）に見られるように，その流れは今，公共交通計画，日々の経路選択，自動走行の制御という異なる対象を通じて，新たな段階に入りつつある．理論を下敷きに新たな駅まち空間を設計しなおすこと．格差を是正しうる公共交通を構想し導入すること．今までにない都市と地域を構想し，計画し，実装し，運営すること．そのための基礎理論を，本書を通じて学んでもらいたい．都市と交通を学ぶとは，変わりゆく現実を理解するだけでなく，失われつつある場所と壊れた風景の前で，それでもなお次の危機に向けた新たな空間を構想し直す力を身につけることなのである．


\addcontentsline{toc}{section}{参考文献}
\bibliographystyle{stone-ship}
\bibliography{../src/bibliography/intro}
